\SourceLeaf{10}
\begin{center}
{\Large\textsc{VORREDE:}}

\medskip
\BookRule{0.07\linewidth}
\end{center}

Die Untersuchungen in diesem Werk gehören zu dem Teil der Mathematik, der sich mit ganzen Zahlen beschäftigt; Brüche bleiben dabei meist, irrationale Größen immer ausgeschlossen. Die sogenannte unbestimmte oder diophantische Analysis, die lehrt, aus den unendlich vielen Lösungen einer unbestimmten Aufgabe diejenigen auszuwählen, die sich durch ganze oder wenigstens rationale Zahlen ausdrücken lassen (meist mit der zusätzlichen Bedingung, dass sie positiv sind), ist nicht diese Wissenschaft selbst, sondern ein sehr besonderer Teil von ihr. Sie steht zu ihr etwa im selben Verhältnis wie die Kunst, Gleichungen umzuformen und zu lösen (Algebra), zur gesamten Analysis. Wie nämlich alle Untersuchungen über die allgemeinen Eigenschaften von Größen zur \emph{Analysis} gehören, so bilden die ganzen Zahlen (und die Brüche, soweit sie durch ganze Zahlen bestimmt werden) den eigentlichen Gegenstand der \textsc{Arithmetik}. Was man aber gewöhnlich unter dem Namen Arithmetik lehrt, geht kaum über die Kunst, Zahlen darzustellen und zu rechnen, hinaus (also Zahlen durch geeignete Zeichen, etwa im Dezimalsystem, darzustellen und die arithmetischen Operationen auszuführen). Hinzu kommen einige Dinge, die entweder gar nicht zur Arithmetik gehören (wie die Lehre von den Logarithmen), oder wenigstens nicht den ganzen Zahlen eigentümlich sind, sondern für alle Größen gelten. Daher scheint es zweckmäßig, zwei Teile der Arithmetik zu unterscheiden: jene Dinge der elementaren Arithmetik zuzuordnen und alle allgemeinen Untersuchungen über die besonderen Eigenschaften ganzer Zahlen der \emph{höheren Arithmetik} vorzubehalten, von der allein hier die Rede sein wird.

Zur höheren Arithmetik gehört, was Euklid im VII. und den folgenden Büchern der Elemente mit der bei den Alten üblichen Eleganz und Strenge dargestellt hat; allerdings beschränkt es sich auf die ersten Anfänge dieser Wissenschaft. Das berühmte Werk des Diophant, das ganz den \SourceLeaf{11}unbestimmten Aufgaben gewidmet ist, enthält viele Fragen, deren Schwierigkeit und feine Kunstgriffe große Achtung vor dem Einfallsreichtum und Scharfsinn des Verfassers wecken, besonders wenn man bedenkt, wie gering die Hilfsmittel waren, die ihm zur Verfügung standen. Diese Aufgaben verlangen aber eher Geschick und gewandte Behandlung als tiefere Grundsätze; außerdem sind sie zu speziell und führen selten zu allgemeineren Schlüssen. Daher scheint dieses Buch eher deshalb einen Wendepunkt in der Geschichte der Mathematik zu bilden, weil es die ersten Spuren einer Zeichensprache und der Algebra zeigt, als weil es die höhere Arithmetik um neue Entdeckungen bereichert hätte. Bei weitem das meiste verdanken wir den Neueren. Unter ihnen haben wenige, aber unsterblich ruhmreiche Männer---P. \textsc{de Fermat}, L. \textsc{Euler}, L. \textsc{La Grange}, A. M. \textsc{Le Gendre} (um einige andere zu übergehen)---den Zugang zum innersten Heiligtum dieser göttlichen Wissenschaft geöffnet und gezeigt, welche reichen Schätze es birgt. Welche Entdeckungen von jedem dieser Mathematiker stammen, will ich hier nicht aufzählen. Man kann dies den Vorreden zu den Zusätzen entnehmen, mit denen der ausgezeichnete La Grange Eulers Algebra bereichert hat, und zu dem kürzlich veröffentlichten Werk des ausgezeichneten Le Gendre, das gleich noch zu nennen ist. Die meisten werden außerdem an den entsprechenden Stellen dieser arithmetischen Untersuchungen angeführt.

Der Zweck dieses Werkes, dessen Veröffentlichung ich schon vor fünf Jahren öffentlich zugesagt hatte, war, Untersuchungen aus der höheren Arithmetik bekannt zu machen, die ich zum Teil vor dieser Zeit, zum Teil danach angestellt habe. Damit niemand sich wundert, dass die Wissenschaft hier beinahe von ihren ersten Anfängen an behandelt wird und viele Untersuchungen wieder aufgenommen werden, mit denen sich andere schon beschäftigt haben, hielt ich folgende Bemerkung für nötig: Als ich mich zu Beginn des Jahres 1795 erstmals dieser Art von Untersuchungen zuwandte, wusste ich nichts von dem, was die Neueren auf diesem Gebiet geleistet hatten, und hatte auch keine Hilfsmittel, durch die ich davon hätte erfahren können. Während einer anderen Arbeit stieß ich zufällig auf eine bemerkenswerte arithmetische Wahrheit (wenn ich mich nicht irre, war es der Satz des Artikels 108). Weil ich sie schon für sich genommen sehr schön fand und vermutete, dass sie mit Größerem zusammenhing, bemühte ich mich mit aller Kraft, ihre Grundlagen zu erkennen und einen strengen Beweis zu finden. Als mir das endlich wie gewünscht gelungen war, hatten mich diese Fragen so sehr gefesselt, dass ich sie nicht mehr verlassen konnte. Da eines immer zum anderen führte, war das meiste von dem, was in den ersten vier Abschnitten dieses Werkes steht, bereits ausgearbeitet, bevor ich irgendetwas von den ähnlichen Arbeiten anderer \SourceLeaf{12}Mathematiker gesehen hatte. Als ich später Gelegenheit bekam, die Schriften dieser großen Geister zu lesen, erkannte ich, dass ich den größten Teil meiner Überlegungen längst erledigten Fragen gewidmet hatte. Doch umso eifriger versuchte ich, ihren Spuren folgend die Arithmetik weiterzuentwickeln; so entstanden verschiedene Untersuchungen, von denen die Abschnitte V, VI und VII einen Teil enthalten. Als ich nach einiger Zeit beschloss, die Früchte meiner nächtlichen Arbeit zu veröffentlichen, ließ ich mich---umso lieber, weil mehrere es wünschten---dazu bewegen, auch von den früheren Untersuchungen nichts zurückzuhalten. Damals gab es kein Buch, aus dem man die Arbeiten anderer Mathematiker über diese Dinge hätte kennenlernen können, die in den Abhandlungen der Akademien verstreut waren; viele meiner Untersuchungen waren ganz neu und die meisten mit neuen Methoden behandelt; schließlich hingen alle sowohl untereinander als auch mit den späteren Untersuchungen so \emph{eng} zusammen, dass auch das Neue ohne eine Wiederaufnahme des Übrigen von Anfang an nicht hinreichend bequem hätte erklärt werden können.

Inzwischen erschien ein ausgezeichnetes Werk eines Mannes, der sich schon zuvor um die höhere Arithmetik sehr verdient gemacht hatte: \emph{Le Gendre, Essai d'une théorie des nombres, Paris, Jahr VI}. Darin hat er nicht nur alles bisher in dieser Wissenschaft Geleistete sorgfältig gesammelt und geordnet, sondern auch sehr vieles Neue aus eigener Arbeit hinzugefügt. Dieses Buch kam erst spät in meine Hände, nachdem der größte Teil meines Werkes schon gedruckt war; daher konnte ich es nicht überall dort erwähnen, wo die Ähnlichkeit der Gegenstände dazu Anlass gegeben hätte. Nur zu wenigen Stellen schien es nötig, in den Zusätzen einige Bemerkungen anzufügen, die dieser überaus höfliche und aufrichtige Mann, wie ich hoffe, freundlich aufnehmen wird.

Während des Drucks dieses Werkes, der mehrfach unterbrochen und durch verschiedene Hindernisse bis in ein viertes Jahr verzögert wurde, habe ich nicht nur schon begonnene Untersuchungen weitergeführt, deren Veröffentlichung ich auf später verschoben hatte, damit das Buch nicht zu umfangreich werde, sondern auch mehrere neue begonnen. Auch mehrere andere, die ich aus demselben Grund nur kurz berührt hatte, weil eine ausführlichere Behandlung weniger nötig schien (zum Beispiel die in den Artikeln 37, 82 und den folgenden sowie an anderen Stellen), wurden später wieder aufgenommen und führten zu allgemeineren Untersuchungen, die mir einer Veröffentlichung sehr wert erscheinen (man vergleiche auch die Bemerkungen in den Zusätzen zu Artikel 306). Da das Buch schließlich, besonders durch die Länge des Abschnitts V, weit umfangreicher wurde als erwartet, musste ich manches ursprünglich dafür Bestimmte \SourceLeaf{13}streichen, darunter den ganzen \emph{achten} Abschnitt (der an verschiedenen Stellen dieses Bandes bereits erwähnt wird und eine allgemeine Behandlung algebraischer Kongruenzen beliebigen Grades enthält). All dies wird leicht einen ebenso großen Band füllen wie den vorliegenden und veröffentlicht werden, sobald sich eine Gelegenheit bietet.

Dass ich bei mehreren schwierigen Fragen synthetische Beweise verwendet und die Analysis, durch die sie gefunden wurden, weggelassen habe, liegt vor allem am Streben nach Kürze, auf die ich so weit wie möglich Rücksicht nehmen musste.

Die in Abschnitt VII behandelte Theorie der Kreisteilung oder der regelmäßigen Vielecke gehört zwar \emph{selbst}, \emph{für sich genommen}, nicht zur Arithmetik; ihre \emph{Grundlagen} müssen aber allein aus der höheren Arithmetik gewonnen werden. Das wird den Mathematikern vielleicht ebenso unerwartet sein, wie ihnen, so hoffe ich, die neuen Wahrheiten willkommen sein werden, die sich aus dieser Quelle gewinnen ließen.

Dies wollte ich dem Leser vorweg erklären. Über die Sache selbst zu urteilen, steht mir nicht zu. Nichts wünsche ich mir mehr, als dass das, was bisherige Lücken füllt oder den Zugang zu Neuem öffnet, denen gefällt, denen der Fortschritt der Wissenschaften am Herzen liegt.

\BookRule{0.16\linewidth}
